\documentclass{beamer}
\setbeamercovered{transparent}


\usetheme{CambridgeUS}
\setbeamercolor{normal text}{bg=gray!20}
\setbeamertemplate{blocks}[default]

\setbeamercolor{block title}{fg=white,bg=red!60!black}
\setbeamercolor{block body}{fg=black,bg=gray!15}


\usepackage[spanish]{babel}

\usepackage[latin1]{inputenc}
\usepackage{amsmath}
\usepackage{amsthm}
\usepackage{amsfonts}
\usepackage{amssymb}
\usepackage{epsfig}
\usepackage{fancyhdr}
\usepackage{psfrag}
\usepackage{pstricks}
\usepackage{colortbl}
\usepackage{longtable}
\usepackage[scriptsize,tight]{subfigure}
\usepackage{multimedia}
\usepackage{color}
\usepackage{fancybox}
\usepackage{comment}
\usepackage{ragged2e}
\usepackage{changepage}
\usepackage{psfrag}


\usefonttheme[onlymath]{serif}


\setbeamertemplate{navigation symbols}{}  

\let\oldtexttt\texttt
\renewcommand{\texttt}[1]{\oldtexttt{\hspace{0pt}#1}} % evita problemas con monospace
\let\oldraggedright\raggedright
\let\raggedright\justifying



\author[Teor�a de Circuitos]{Teor�a de Circuitos}
\title[Tema 1]{\scshape Tema 1}
\subtitle{Conceptos fundamentales y leyes de Kirchhoff}

\date[]{}
\institute[ETSI]{Departamento de Ingenier�a El�ctrica \\ Escuela T�cnica Superior de Ingenier�a\\ Universidad de Sevilla}


\begin{document}


\begin{frame}
\titlepage
\end{frame}


\section{Magnitudes el�ctricas}


\begin{frame}
\frametitle{Carga el�ctrica}
\begin{columns}
\begin{column}{10cm}
\structure{\textbf{Carga el�ctrica}}, $q(t)$: propiedad de las part�culas elementales que constituyen la materia y que se manifiesta por medio de fuerzas el�ctricas.

\medskip

\begin{itemize}
	\item El electr�n, e$^-$, constituye la carga el�ctrica m�s peque�a posible, representada como $q_e$.
	
	\item Se toma la carga del electr�n como negativa, y positiva la del prot�n.
	
	\item Aplicando energ�a se puede conseguir que los electrones se desplacen.
\end{itemize}

\medskip

\structure{Unidad en el SI:} culombio [C] que equivale a $6.242\cdot 10^{18}$\,$q_e$.
\end{column}
\begin{column}{2cm}
\begin{figure}[hbt]
\vspace*{-0.5cm}
\epsfig{file=coulomb.eps,height=2cm}
\label{coulomb}
\end{figure}
\vspace*{-0.5cm}
\begin{center}
\tiny{Charles-Augustin de Coulomb  \\
(1736-1806)}
\end{center}
\end{column}
\end{columns}


\end{frame}




\end{document}
















